\documentclass[10pt]{mypackage}

\usepackage{mlmodern}
%\usepackage{newpxtext,eulerpx,eucal}
%\renewcommand*{\mathbb}[1]{\varmathbb{#1}}

\usepackage{homework}
%\usepackage{notes}

%\usepackage[ backend=bibtex, style = alphabetic, sorting=ynt ]{biblatex}
%\addbibresource{  }

\usepackage{parskip}

\fancyhf{}
\fancyhead[R]{Avinash Iyer}
\fancyhead[L]{Algebra II: Homework 7}
\fancyfoot[C]{\thepage}

\setcounter{secnumdepth}{0}

\begin{document}
\RaggedRight
\begin{problem}[Problem 2]
  Let $\left( \rho,V \right)$ be a representation of a group $G$. Recall that the dual representation $\left( \rho^{\ast},V^{\ast} \right)$ is defined by $\rho^{\ast}\left( g \right)f = f\circ \rho\left( g^{-1} \right)$ for all $f\in V^{\ast}$. Prove that
  \begin{enumerate}[(i)]
    \item $\left( \rho^{\ast},V^{\ast} \right)$ is a representation;
    \item if $\beta$ is a basis for $V$, and $\beta^{\ast}$ is the dual basis for $V^{\ast}$, then $\left[ \rho^{\ast}\left( g \right) \right]_{\beta^{\ast}} = \left[ \rho\left( g^{-1} \right) \right]_{\beta}^{T}$. 
  \end{enumerate}
\end{problem}
\begin{solution}\hfill
  \begin{enumerate}[(i)]
    \item If $g,h\in G$, and $f\in V^{\ast}$ is arbitrary, then
      \begin{align*}
        \rho^{\ast}\left( gh \right)f &= f\circ \rho\left( \left( gh \right)^{-1} \right)\\
                                      &= f\circ \rho\left( h^{-1}g^{-1} \right)\\
                                      &= \left( f\circ \rho\left( h^{-1} \right) \right)\circ \rho\left( g^{-1} \right)\\
                                      &= \rho^{\ast}\left( g \right)\left( f\circ \rho\left( h^{-1} \right) \right)\\
                                      &= \rho^{\ast}\left( g \right)\rho^{\ast}\left( h \right)f.
      \end{align*}
      Therefore, we have $\rho^{\ast}\left( gh \right) = \rho^{\ast}\left( g \right)\rho^{\ast}\left( h \right)$ for all $g,h\in G$.
    \item Our goal is to show that $\left[ A^{\ast} \right]_{\beta^{\ast}} = \left[ A \right]_{\beta}^{T}$. We note that if $\left[ A \right]_{\beta}$ is a matrix on some basis $\beta$, then the matrix coefficient $a_{ij}$ is extracted by evaluating $v_i^{\ast}\left( Av_j \right)$, where $v_i^{\ast}$ is the dual basis element for $v_i$. In particular, this means that, if $b_{ij}$ denotes the $\left( i,j \right)$ entry of $\left[ A^{\ast} \right]_{\beta^{\ast}}$, we have $b_{ij} = \hat{v}_i\left( A^{\ast}v_j^{\ast} \right)$, where $\hat{v}_i$ denotes the double dual element for $v_i$. Yet, this gives
      \begin{align*}
        b_{ij} &= \hat{v}_i\left( A^{\ast}v_j^{\ast} \right)\\
               &= \left( A^{\ast}v_j^{\ast} \right)\left( v_i \right)\\
               &= v_j^{\ast}\left( Av_i \right)\\
               &= a_{ji}.
      \end{align*}
      Thus, $\left[ A^{\ast} \right]_{\beta^{\ast}} = \left[ A \right]_{\beta}^{T}$.
  \end{enumerate}
\end{solution}
\begin{problem}[Problem 3]
  Describe all irreducible complex representations of $S_4$ and compute its character table.
\end{problem}
\begin{solution}
  The following partitions of $4$ give that there are $5$ conjugacy classes:
  \begin{enumerate}[(i)]
    \item $4 = 1 + 1 + 1 + 1$ (the trivial cycle type);
    \item $4 = 2 + 1 + 1$ (transpositions);
    \item $4 = 2 + 2$ (the Klein $4$-group, $V_4$);
    \item $4 = 3 + 1$ (the $3$-cycles);
    \item $4 = 4$ (the $4$-cycles).
  \end{enumerate}
  We will label these classes by $K_1,\dots,K_5$. We already know of three representations for $S_4$:
  \begin{enumerate}[(i)]
    \item $\rho_{t}$, which is the trivial representation (which has dimension 1);
    \item $\rho_{\sgn}$, the sign representation (which has dimension 1);
    \item $\rho_{\operatorname{std}}$, the standard representation (which has dimension 3).
  \end{enumerate}
  Since $S_4/V_4\cong S_3$, we know that these three representations also define the three irreducible representations for $S_3$; using Problem 1, we may then find a representation which we will call $\rho_{q}$ (quotient) for $S_4$ by taking $\rho(g) = \rho_{\operatorname{std}}\left(gV_4\right)$. This representation has dimension 2.

  Finally, we know that the tensor product of an irreducible representation by a one-dimensional representation is also irreducible, so taking the tensor product of $\rho_{\operatorname{std}}$ with $\rho_{\operatorname{sgn}}$ gives the representation $\rho_{\operatorname{twist}} = \rho_{\operatorname{std}}\otimes \rho_{\operatorname{sgn}}$. We thus get the following character table.
  \begin{center}
    \begin{tabular}{c|c|c|c|c|c}
      & $K_1$ & $K_2$ & $K_3$ & $K_4$ & $K_5$\\
      \hline
      $\chi_{t}$ & $ 1 $ & $1$ & $1$ & $1$ & $1$ \\
      $\chi_{\operatorname{sgn}}$ & $1$ & $-1$ & $1$ & $1$ & $-1$ \\
      $\chi_{q}$ & $2$ & & $1$ & & \\
      $\chi_{\operatorname{std}}$ & $3$ & & & & \\
      $\chi_{\operatorname{twist}}$ & $3$ & & & & \\
      \hline
    \end{tabular}
  \end{center}
\end{solution}
\end{document}
